\begin{titlepage}
\thispagestyle{empty}
\begin{tikzpicture}[remember picture,overlay]
  \fill[indigo] (current page.north west) rectangle ([yshift=-10.6cm]current page.north east);
  \begin{scope}
    \clip (current page.north west) rectangle ([yshift=-10.6cm]current page.north east);
    \draw[white,opacity=0.05,step=0.42cm]
      ([yshift=-10.6cm]current page.north west) grid (current page.north east);
    \begin{scope}[shift={([xshift=14.5cm,yshift=-5.4cm]current page.north west)}]
      \node[text=white,opacity=0.92,font=\large] at (0,1.45) {$Z=\mathrm{Tr}\!\left(T^{N}\right)$};
      \draw[gold,line width=1.1pt,rounded corners=2pt] (-1.6,-0.68) rectangle (1.6,0.85);
      \node[text=gold,font=\small] at (0,0.08)
        {$\begin{pmatrix} e^{\beta(J+h)} & e^{-\beta J}\\[2pt] e^{-\beta J} & e^{\beta(J-h)}\end{pmatrix}$};
      \draw[white,opacity=0.6,line width=1pt,-{Stealth[length=5pt]}] (0,-0.84)--(0,-1.45);
      \node[text=rose,font=\normalsize] at (0,-1.93) {$\lambda_+^{\,N}+\lambda_-^{\,N}$};
      \node[text=white,opacity=0.5,font=\scriptsize,align=center] at (0,-2.62)
        {$2^{N}$ configurations\\collapse to two numbers};
    \end{scope}
  \end{scope}

  \node[anchor=north west,text=white,align=left,inner sep=0pt]
    at ([xshift=2.05cm,yshift=-1.75cm]current page.north west)
    {\fontsize{10}{12}\selectfont\color{gold}\textsc{part x $\cdot$ version 2.0 $\cdot$ supersedes v1.0}\\[2pt]
     \fontsize{9}{11}\selectfont\color{white!70}a corrected audit of all four libraries};

  \node[anchor=north west,text=white,align=left,inner sep=0pt,text width=11.4cm]
    at ([xshift=2.05cm,yshift=-3.0cm]current page.north west)
    {\fontsize{42}{46}\selectfont\bfseries Learning $Z$\\[10pt]
     \fontsize{13.5}{18}\selectfont\color{white!92}
     The partition function --- and the discovery\\that you already own almost every book\\[10pt]
     \fontsize{10}{14}\selectfont\color{white!76}\normalfont
     A re-audit across \textbf{Pure\_Maths\_TC\_SelfStudy}, \textbf{TensoQi101},\\
     \textbf{QBIT-HIS} and \textbf{qm-qi-qml} --- 681 books, every claim\\
     verified by opening the file.};

  \foreach \i/\c/\t/\n/\s in {0/sand/{Pure\_Maths}/153/{measure, analysis,\\Feynman's stat mech},
                              1/tenso/{TensoQi101}/121/{QM, quantum info,\\DFT, Shankar},
                              2/qbit/{QBIT-HIS}/66/{quantum stat mech,\\stat mech of the cell},
                              3/slate/{qm-qi-qml}/341/{QML, MIT 8.04\\(little new for $Z$)}}{
    \node[anchor=north west,inner sep=0pt]
      at ([xshift={2.05cm+\i*4.28cm},yshift=-11.85cm]current page.north west)
      {\begin{tikzpicture}
         \node[fill=\c,text=white,rounded corners=4pt,minimum width=3.95cm,
               minimum height=2.3cm,align=left,inner sep=9pt,anchor=north west,
               font=\scriptsize] at (0,0)
           {\textbf{\footnotesize \t}\\[3pt]{\fontsize{16}{17}\selectfont\bfseries \n}\; books\\[4pt]\s};
       \end{tikzpicture}};}

  \node[anchor=north west,inner sep=0pt]
    at ([xshift=2.05cm,yshift=-14.95cm]current page.north west)
    {\begin{tikzpicture}
       \node[draw=alert,line width=1.3pt,rounded corners=3pt,fill=alert!6,
             text width=16.35cm,align=left,inner sep=13pt,font=\small,anchor=north west] at (0,0)
         {\textbf{\color{alert}What version 1.0 got wrong.}\;
          v1.0 reported that Modules 5--8 of \bk{Pure\_Maths\_TC\_SelfStudy} were empty, and sent you
          shopping for five books.\\[5pt]
          {\footnotesize\color{ink}Four of those five were already on your shelf, dated \textbf{March
          2024} --- eighteen months before the audit. \textbf{Feynman's \textit{Statistical
          Mechanics}} was sitting in \bk{Module\_8}. The corrected finding is in §1, and the gap list
          has collapsed from five books to \textbf{two}.}};
     \end{tikzpicture}};

  \node[anchor=north west,inner sep=0pt]
    at ([xshift=2.05cm,yshift=-18.4cm]current page.north west)
    {\begin{tikzpicture}
       \node[draw=gold,line width=1.1pt,rounded corners=3pt,fill=gold!7,
             text width=16.35cm,align=left,inner sep=12pt,font=\small,anchor=north west] at (0,0)
         {\textbf{\color{ink}The standout find.}\;
          \textbf{Shankar §21.2} (\textit{Principles of QM}, p.~627, in \bk{TensoQi101}) solves the 1D
          Ising partition function \emph{and} maps it onto the quantum path integral. \textbf{Feynman
          ch.~1, 3 and 5} does the same arc from the other direction.\\[5pt]
          {\footnotesize\color{slate}You own two independent treatments of the exact calculation in
          §5 of this dossier --- one from quantum mechanics, one from statistical mechanics.}};
     \end{tikzpicture}};

  \node[anchor=south west,align=left,text=slate,font=\scriptsize,inner sep=0pt,text width=16.9cm]
    at ([xshift=2.05cm,yshift=2.05cm]current page.south west)
    {\color{indigo}\rule{16.9cm}{1.2pt}\\[6pt]
     Every book cited was opened and its page count recorded. Numerical results in §5--§6 are
     reproducible from \bk{Partition\_Function\_01\_Ising\_Transfer\_Matrix.ipynb}, unchanged from v1.0.};
\end{tikzpicture}
\end{titlepage}
